\section{Closure: mass spectrum as resolution depth and protocol cost}
\label{sec:mass_spectrum_closure}

This section closes the interface by connecting the field-level labeling map (Section~\ref{sec:sm_labeling_closure}) to a concrete, reproducible mass-spectrum template.
We work with dimensionless ratios relative to the electron reference, and we phrase running/threshold effects in the Fibonacci resolution coordinate of the golden branch.

\subsection{Fibonacci resolution coordinate}
\label{subsec:mass_resolution_coordinate}

Fix the reference scale $\mu_0=m_e$ and define the resolution coordinate
\begin{equation}\label{eq:r_of_mu_z128}
r(\mu):=\frac{\log(\mu/m_e)}{\log\varphi},
\end{equation}
where $\varphi=(1+\sqrt{5})/2$ is the golden ratio.
Equivalently, the exponential map is
\begin{equation}\label{eq:mu_of_r_z128}
\mu(r)=m_e\,\varphi^{\,r}.
\end{equation}
This is the same resolution-flow dictionary recorded in the constants-geometry sector \cite{Ma2025PCG}.

\subsection{A closed depth assignment from stable-type invariants}
\label{subsec:mass_depth_assignment}

Let $w\in X_6$ and recall the intrinsic invariants $(V(w),g(w),D_\pi(w))$ (Definition~\ref{def:x6_invariants}).
We use the effective protocol depth $r_\ast(w)$ from Definition~\ref{def:r_star}.
The interpretation is that folding degeneracy counts extra protocol effort (additional scan density) in units of one spatial addressing order.

\begin{definition}[Closed mass template]\label{def:mass_template}
Fix the electron reference field $e:=e_R^{(1)}$ and let $w_e\in X_6$ denote its stable-type label under $\mathcal{L}_{\mathrm{SM}}$.
Define the normalized depth
\begin{equation}\label{eq:r_hat_def}
\widehat r(f):=\kappa\bigl(r_\ast(w_f)-r_\ast(w_e)\bigr),
\end{equation}
where $w_f\in X_6$ satisfies $\mathcal{L}_{\mathrm{SM}}(w_f)=f$ and $\kappa:=m/n=2$ for the minimal balanced coupling $(m,n)=(6,3)$.
Then define the closed template mass prediction by
\begin{equation}\label{eq:mass_template}
\mu_{\mathrm{pred}}(f):=m_e\,\varphi^{\,\widehat r(f)}.
\end{equation}
\end{definition}

Equation~\eqref{eq:mass_template} yields a discrete spectrum of reference scales.
Threshold and scheme effects are recorded as multiplicative matching factors (equivalently, additive shifts in $r$) as in effective field theory \cite{PeskinSchroeder1995,PDGRPP2024}.

\subsection{Predicted spectrum and PDG/CODATA comparisons}
\label{subsec:mass_spectrum_table}

Table~\ref{tab:mass_spectrum} records the closed template values together with standard reference values (PDG/CODATA conventions).
For quarks we use a scheme-dependent reference (e.g.\ $\overline{\mathrm{MS}}$ running masses) and treat scheme dependence as a matching input rather than an integer-depth anchor, consistent with the resolution-map calibration philosophy \cite{Ma2025PCG}.
For neutrinos we record an order-of-magnitude reference scale inferred from oscillation data, but we do not fix a unique absolute mass prediction at this minimal resolution.
For $W$, $Z$, and $H$ we include the canonical electroweak anchor depths as discrete reference points in the $r$-coordinate.

\begin{table}[t]
\centering
\small
\setlength{\tabcolsep}{6pt}
\renewcommand{\arraystretch}{1.15}
\begin{tabular}{lrrrrr}
\toprule
field & reference $\mu$ [GeV] & $r(\mu)$ & $\widehat r$ & $\mu_{\mathrm{pred}}$ [GeV] & $\mu/\mu_{\mathrm{pred}}$ \\
\midrule
$\nu$ (scale) & $5\times 10^{-11}$ & -33.540 & $-$ & $-$ & $-$ \\
$e$ & $5.10999\times 10^{-4}$ & 0.000 & 0 & $5.10999\times 10^{-4}$ & $1$ \\
$u$ & $0.00216$ & 2.996 & 6 & $0.009169504105$ & $0.235563$ \\
$d$ & $0.00467$ & 4.598 & 4 & $0.003502438908$ & $1.33336$ \\
$s$ & $0.093$ & 10.814 & 10 & $0.06284871611$ & $1.47974$ \\
$\mu$ & $0.10565838$ & 11.080 & 10 & $0.06284871611$ & $1.68115$ \\
$c$ & $1.27$ & 16.247 & 12 & $0.1645400749$ & $7.71848$ \\
$\tau$ & $1.77686$ & 16.945 & 14 & $0.4307715087$ & $4.12483$ \\
$b$ & $4.18$ & 18.722 & 20 & $7.729881083$ & $0.540759$ \\
$W$ & $80.377$ & 24.866 & 25 & $85.72569485$ & $0.937607$ \\
$Z$ & $91.1876$ & 25.128 & 25 & $85.72569485$ & $1.06371$ \\
$H$ & $125.25$ & 25.788 & 26 & $138.707088$ & $0.902982$ \\
$t$ & $172.76$ & 26.456 & 24 & $52.98139313$ & $3.26077$ \\
%
\\
\bottomrule
\end{tabular}
\caption{Mass-spectrum closure in the resolution-depth language. The predicted values use the normalized depth \eqref{eq:r_hat_def} built from the stable-type cost \eqref{eq:r_star_def}, together with the exponential map \eqref{eq:mu_of_r_z128}. The ratio $\mu/\mu_{\mathrm{pred}}$ records the multiplicative matching input at the stated reference conventions. All rows are generated by a reproducible script.}
\label{tab:mass_spectrum}
\end{table}


